% 透かしの設定
\DraftwatermarkOptions{angle={0}}
\DraftwatermarkOptions{scale={0.15}}
\DraftwatermarkOptions{hpos={0.90\paperwidth}}
\DraftwatermarkOptions{vpos={0.20\paperwidth}}
\DraftwatermarkOptions{color={[rgb]{1, 0, 0}}}
\DraftwatermarkOptions{text={%
  \includegraphics{./../../images/maruhi_mark.png} \\ サンプル
}}

% ヘッダとフッタを設定
\renewcommand{\sectionmark}[1]{\markright{\thesection.\ #1}}

\fancypagestyle{normal}{
  \fancyhf{}
  \fancyhead[LE]{\rightmark}
  \fancyhead[RO]{\rightmark}
  \fancyfoot[C]{\thepage}
  \renewcommand{\headrulewidth}{0.4pt}
  \renewcommand{\footrulewidth}{0.4pt}
}

% chapter の最初のページにスタイルを適用
% https://github.com/texjporg/uplatex/blob/master/ujclasses.dtx#L2895
% arg1: 対象のコマンド
% arg2: 置換前のコード
% arg3: 置換後のコード
% arg4: デバック用
% arg5: デバック用
\makeatletter
\patchcmd{\chapter}{\thispagestyle{jpl@in}}{\thispagestyle{normal}}{}{}
\makeatother

% 表目次と図目次を追加
\makeatletter
\renewcommand{\sphinxtableofcontents}{
  \pagenumbering{roman}
  \begingroup
    \parskip
    \z@skip
    \sphinxtableofcontentshook
    \tableofcontents
    \listoffigures
    \listoftables
  \endgroup
  \if@openright
    \cleardoublepage
  \else
    \clearpage
  \fi
  \pagenumbering{arabic}
}
\makeatother

% 目次にスタイルを適用
\AtBeginDocument{\addtocontents{toc}{\protect\thispagestyle{normal}}}
\AtBeginDocument{\addtocontents{lof}{\protect\thispagestyle{normal}}}
\AtBeginDocument{\addtocontents{lot}{\protect\thispagestyle{normal}}}

% 関連図書を参照一覧に変更
\renewcommand{\bibname}{参照一覧}
